
% ════════════════════════════════════════════
% 模型的评价（07.模型评价.tex）—— 写作提示词见本文件末尾【本章节写作提示词】。
% ════════════════════════════════════════════
\section{模型的评价}

\subsection{模型的优点}
\begin{itemize}[itemindent=2em]
\item \textbf{优点1：} 可观测性判定模型将车辆几何、视线遮挡与云台俯仰统一为三个可独立计算的约束，几何意义清晰，适用于任意轴对齐建筑布局。
\item \textbf{优点2：} 交会定位精度模型给出了位置误差均方根的闭式表达，直接揭示基线长度、交会角与目标距离对精度的作用规律，便于工程应用。
\item \textbf{优点3：} 并发调度与重调度模型采用穷举搜索，在小规模任务下保证全局最优，并给出完整的车辆时间轴，可解释性强。
\item \textbf{优点4：} 误差传播模型解析与蒙特卡洛仿真相互印证，误差来源方差分解直观，为车辆部署提供了量化约束。
\end{itemize}

\subsection{模型的缺点}
\begin{itemize}[itemindent=2em]
\item \textbf{缺点1：} 调度模型采用固定候选观测站与离散网格，最优性受网格分辨率影响，且未同时优化定位精度与响应时间的多目标权衡。
\item \textbf{缺点2：} 重调度模型假设车辆位置在触发时刻可由基线调度精确插值，未考虑行驶过程的随机波动与多车协同中的避碰细节。
\end{itemize}
